% ============================================================
% Chapter 6: Experimental Design
%   Section: Evaluation Harness
%     Subsection: Metrics and Confidence Intervals
% Included from chapters/06-experimental-design/04-evaluation-harness/evaluation-harness.tex
% Source: docs/decisions/evaluation.md, "What is reported apart",
% "Metrics"; metrics.py, report.py.
% ============================================================

\subsection{Metrics and Confidence Intervals}
\label{subsec:harness-metrics}

\paragraph{What is reported apart.} Each claim gets one outcome, the
first that applies: \emph{error} (the harness or the pipeline crashed),
\emph{search unavailable} (no web query was answered),
\emph{model unreachable}, or \emph{scored}. Only scored claims enter
the metrics. The others are listed by identifier and never counted as
errors of the system: on 25~September, 68 of 69 searches came back
empty, and folded into the error rate that would have measured the
search engine's uptime. On top of the outcome, two flags:

\begin{itemize}
    \item \textbf{Temporal leak}: a cited source was published after
    the claim (the annotation guide's Rule~3). Undetermined when the
    claim or the source has no date.
    \item \textbf{Verdict leak} (X-Fact only): a ranked source is on the
    fact-checker's own site, so the verdict may have been read rather
    than reached.
\end{itemize}

Every metric is reported on all scored claims and again without the
flagged ones. An answer the model gave in no usable format is scored as
the \texttt{UNVERIFIED} it became and counted per model.

\paragraph{Verdict metrics.} Over the five verdicts:

\begin{itemize}
    \item \textbf{accuracy}, and \textbf{macro-averaged $F_1$} over the
    classes present in the gold labels --- a class with no gold support
    is shown with its predictions but not averaged in --- with
    per-class precision, recall, $F_1$, support and predicted count, and
    the confusion matrix;
    \item \textbf{Cohen's $\kappa$} against the gold labels
    \cite{cohen1960coefficient}, the same statistic as the annotator's
    self-agreement, so that the system and the annotator are compared
    on one scale;
    \item \textbf{coverage} --- the share of claims given a definitive
    verdict rather than \texttt{UNVERIFIED} --- and \textbf{selective
    accuracy}, the accuracy on those claims, the standard pair for a
    classifier allowed to abstain \cite{geifman2017selective}.
\end{itemize}

\paragraph{Intervals.} Every statistic comes with a percentile 95\%
bootstrap interval over claims \cite{efron1993bootstrap}: 10,000
resamples, seeded (2026 by default) and with the seed in the report,
records ordered by identifier first so that the same run gives the same
interval however its claims finished. Two runs are compared with a
\emph{paired} bootstrap of their difference over the claims both
scored, alongside how many claims only one of them got right; a run
compared with itself gives exactly zero, which a test holds.

\paragraph{Breakdowns.} By language for both sets; by topic group,
claim type, source tier and origin (natural or synthetic) for the
custom set; by fact-checking site for X-Fact. Each with its size,
accuracy and interval, macro-$F_1$, coverage and selective accuracy.

\paragraph{Cost.} Each record carries its model calls, tokens and the
provider's answer time, measured inside the production client, so the
report gives the median and 90th-percentile latency per claim, tokens
per claim and per second, and the price per hundred claims from the
price table --- ``n/a'' without a price, never zero.
